\documentclass[11pt,a4paper]{article}

\usepackage[T1]{fontenc}
\usepackage[utf8]{inputenc}
\usepackage{mathptmx}
\usepackage[scaled=0.9]{helvet}
% --- localización manual (babel-spanish no disponible en esta distribución)
\renewcommand{\contentsname}{Índice}
\renewcommand{\appendixname}{Apéndice}
\renewcommand{\figurename}{Figura}
\renewcommand{\tablename}{Cuadro}
\renewcommand{\refname}{Referencias}
\usepackage[margin=2.4cm]{geometry}
\usepackage{amsmath,amssymb,amsthm,mathtools}
\usepackage{booktabs}
\usepackage{tikz}
\usetikzlibrary{arrows.meta,calc,positioning,patterns,decorations.pathreplacing}
\usepackage{microtype}
\usepackage{enumitem}
\usepackage{tcolorbox}
\usepackage[hidelinks]{hyperref}
\usepackage{fancyhdr}
\usepackage{caption}

\tcbuselibrary{skins,breakable}

\captionsetup{font=small,labelfont=bf}

\pagestyle{fancy}
\fancyhf{}
\fancyhead[L]{\footnotesize Hoja de medición del rover · Vision Rover Challenge}
\fancyhead[R]{\footnotesize\thepage}
\renewcommand{\headrulewidth}{0.4pt}

% ---------------------------------------------------------------- entornos
\newtheorem{teorema}{Teorema}[section]
\newtheorem{proposicion}[teorema]{Proposición}
\newtheorem{lema}[teorema]{Lema}
\newtheorem{corolario}[teorema]{Corolario}
\theoremstyle{definition}
\newtheorem{definicion}[teorema]{Definición}
\newtheorem{observacion}[teorema]{Observación}
\newtheorem{ejemplo}[teorema]{Ejemplo}

\newenvironment{demo}{\par\noindent\textit{Demostración.}\;}{\hfill$\blacksquare$\par\medskip}

\newtcolorbox{resultado}[1][]{
  enhanced, breakable, colback=black!3, colframe=black!55,
  boxrule=0.5pt, arc=1pt, left=6pt, right=6pt, top=5pt, bottom=5pt, #1}

% ---------------------------------------------------------------- macros
\newcommand{\R}{\mathbb{R}}
\newcommand{\Rot}{\mathbf{R}}
\newcommand{\uu}{\mathbf{u}}
\newcommand{\vv}{\mathbf{v}}
\newcommand{\nn}{\mathbf{n}}
\newcommand{\dd}{\mathbf{d}}
\newcommand{\cc}{\mathbf{c}}
\newcommand{\pp}{\mathbf{p}}
\newcommand{\qq}{\mathbf{q}}
\newcommand{\NN}{\mathbf{N}}
\newcommand{\ww}{\mathbf{w}}
\newcommand{\ee}{\mathbf{e}}
\newcommand{\gv}{\mathbf{g}}
\DeclareMathOperator*{\argmin}{arg\,min}
\DeclareMathOperator*{\argmax}{arg\,max}
\newcommand{\aaa}{\mathbf{a}}
\newcommand{\cross}[2]{#1 \wedge #2}
\newcommand{\dif}{\operatorname{d}\!}
\newcommand{\atantwo}{\operatorname{atan2}}
\newcommand{\conv}{\operatorname{conv}}
\newcommand{\intr}{\operatorname{int}}
\newcommand{\diam}{\operatorname{diam}}
\newcommand{\sgn}{\operatorname{sgn}}
\newcommand{\celda}{\mathrm{c}}
\newcommand{\grados}{^{\circ}}

\newtcolorbox{medicion}[2]{
  enhanced, breakable, colback=white, colframe=black!70,
  boxrule=0.8pt, arc=1pt, left=7pt, right=7pt, top=6pt, bottom=6pt,
  title={\bfseries #1\ \ ---\ \ #2}, fonttitle=\normalsize,
  colbacktitle=black!70, coltitle=white}

\newcommand{\campo}[1]{\par\smallskip\noindent\textbf{#1.}\ }
\newcommand{\anota}[1]{\par\smallskip\noindent\fbox{\parbox[c][1.1cm][c]{0.97\linewidth}{\footnotesize\textsc{medido:}\hfill #1}}}

\begin{document}

\begin{center}
{\Large\bfseries Hoja de medición del rover CenfoBot}\\[0.25cm]
{\large Verificación experimental del modelo geométrico}\\[0.5cm]
\rule{\linewidth}{1pt}
\end{center}

\vspace{0.2cm}

\noindent
Esta hoja está pensada para llenarse \textbf{con el rover en la mano y una regla}.
Cada medición trae el valor que el modelo predice y qué significa si el número real
no coincide. El orden no es arbitrario: las tres primeras deciden si el hallazgo H-11
---el que dice que conocer el ángulo del cubo es obligatorio--- se sostiene o se cae.

\begin{tcolorbox}[colback=black!5,colframe=black!65,boxrule=0.7pt,arc=1pt]
\textbf{Cómo medir con regla lo que parece necesitar calibrador.} El modelo dice
$W_u = 93{,}50\ \mathrm{mm}$ y la pregunta de fondo es si $W_u$ supera
$94{,}85\ \mathrm{mm}$. Distinguir $93{,}5$ de $94{,}9$ sobre un tramo de $94$ es
imposible con regla. Pero esa misma diferencia, medida sobre un \textbf{hueco de
$9\ \mathrm{mm}$}, se ve a simple vista.

\smallskip
Todas las mediciones críticas de esta hoja están planteadas como \emph{huecos}, no como
distancias grandes. Es la misma información, leída donde la regla sí alcanza.
\end{tcolorbox}

\section*{Bloque A --- Las tres que deciden H-11}
\addcontentsline{toc}{section}{Bloque A}

\begin{medicion}{M1}{La prueba de los 10 milímetros \textnormal{(la más importante de todas)}}
\campo{Qué hacer} Meté un cubo entre las paletas, \textbf{girado a $45\grados$} respecto
del eje del rover (en rombo, con una punta mirando al frente). Empujalo lateralmente
hasta que toque una paleta. Medí el \textbf{hueco que queda del otro lado}.

\campo{Qué predice el modelo} $8{,}65\ \mathrm{mm}$.

\campo{El umbral que importa} $\mathbf{10{,}00\ \mathrm{mm}}$. No hace falta leer el
número con precisión: alcanza con saber de qué lado del umbral cae.
\begin{itemize}[leftmargin=1.6em,itemsep=2pt,topsep=3pt]
\item \textbf{Si el hueco es menor que $10\ \mathrm{mm}$} $\to$ H-11 confirmado. El
      presupuesto de control es negativo en el peor caso y el ángulo del cubo es
      obligatorio.
\item \textbf{Si es mayor que $10\ \mathrm{mm}$} $\to$ H-11 se suaviza. Sigue siendo
      apretado, pero deja de ser imposible.
\end{itemize}

\campo{Truco pasa/no pasa} Si conseguís algo de $10\ \mathrm{mm}$ de espesor conocido
(dos monedas apiladas medidas antes, un taco de acrílico de $3+3+3$, un lápiz de sección
conocida), intentá meterlo en el hueco junto al cubo. \textbf{Si no entra, H-11 está
confirmado} y no hiciste ninguna lectura de regla.

\campo{Por qué esta medición} $e_{\max}$ en el peor caso es exactamente
\emph{la mitad de este hueco}. Es la única magnitud del modelo que se mide de forma
directa, sin restar dos números grandes.
\anota{$\rule{4cm}{0.4pt}$ mm \quad$\Rightarrow$\quad $e_{\max} = \rule{2cm}{0.4pt}$ mm}
\end{medicion}

\begin{medicion}{M2}{Hueco con el cubo alineado \textnormal{(control de la medición anterior)}}
\campo{Qué hacer} Lo mismo que M1 pero con el cubo \textbf{alineado} (caras paralelas a
las paletas), empujado contra una paleta. Medir el hueco del otro lado.

\campo{Qué predice el modelo} $33{,}50\ \mathrm{mm}$.

\campo{El control} La resta $\text{M2} - \text{M1}$ debe dar
$L(\sqrt2 - 1) = \mathbf{24{,}85\ \mathrm{mm}}$, \emph{cualquiera sea el valor de
$W_u$}. Es una identidad geométrica que solo depende de la arista del cubo.

\smallskip
Si la resta da $24{,}85 \pm 1$, las dos mediciones son buenas y se pueden creer. Si da
otra cosa, algo se midió mal ---o el cubo no mide $60\ \mathrm{mm}$, o las paletas no
son paralelas--- y hay que repetir antes de seguir.
\anota{$\rule{4cm}{0.4pt}$ mm \quad$\Rightarrow$\quad M2 $-$ M1 $= \rule{2cm}{0.4pt}$ mm}
\end{medicion}

\begin{medicion}{M3}{¿La lengüeta queda enrasada?}
\campo{Qué hacer} Mirá por fuera dónde el chasis (la plataforma horizontal) atraviesa el
panel lateral. \textbf{Pasá la uña} por encima de la junta, sin medir nada.

\campo{Las tres respuestas posibles}
\begin{itemize}[leftmargin=1.6em,itemsep=2pt,topsep=3pt]
\item \textbf{Enrasada} (la uña no salta) $\to$ junta pasante. $W_u = W_{\mathrm{ext}} -
      6{,}00$, y todo el modelo actual es correcto.
\item \textbf{Sobresale} $\to$ la lengüeta es más larga que el espesor. $W_u$ sigue
      siendo el cuerpo del chasis, pero $W_{\mathrm{ext}}$ es mayor y hay que corregir
      $R_{\mathrm{rov}}$.
\item \textbf{Hundida} $\to$ la lengüeta no atraviesa del todo. Entonces $W_u$ es
      \emph{mayor} que $93{,}50$ y el presupuesto mejora.
\end{itemize}

\campo{Por qué} Esta observación cualitativa, que no requiere ningún instrumento,
es la que justifica toda la deducción de $W_u$ a partir del archivo de fabricación.
\anota{enrasada $\square$ \qquad sobresale $\square$ \qquad hundida $\square$}
\end{medicion}

\section*{Bloque B --- El sesgo escondido}
\addcontentsline{toc}{section}{Bloque B}

\begin{tcolorbox}[colback=black!5,colframe=black!65,boxrule=0.7pt,arc=1pt]
\textbf{Por qué este bloque puede ser peor que H-11.} La visión no ve al rover: ve su
marcador ArUco, y reporta \emph{la posición del marcador}. Si el marcador no está en el
centro geométrico del rover, cada posición publicada trae un \textbf{sesgo fijo}, y ese
sesgo se resta directamente del presupuesto lateral.

\smallskip
Con $4{,}32\ \mathrm{mm}$ de tolerancia y $5\ \mathrm{mm}$ ya gastados en ruido, un
marcador descentrado $3\ \mathrm{mm}$ no empeora las cosas: las vuelve imposibles
también para los cubos alineados. Y a diferencia del ruido, un sesgo \textbf{se puede
corregir en software} en cuanto se conoce el número.
\end{tcolorbox}

\begin{medicion}{M4}{Descentrado lateral del marcador ArUco}
\campo{Qué hacer} Con el rover visto desde arriba, medir del \textbf{centro del
marcador} a la cara exterior del panel izquierdo, y del centro del marcador a la cara
exterior del panel derecho.

\campo{Qué predice el modelo} Las dos iguales, $W_{\mathrm{ext}}/2 = 49{,}75\ \mathrm{mm}$.

\campo{Qué significa} La mitad de la diferencia entre ambas es el sesgo lateral $b$, y
el presupuesto real pasa a ser $e_{\max} - \sigma_{\mathrm{vis}} - b$. Cualquier valor
de $b$ por encima de $1\ \mathrm{mm}$ hay que compensarlo en el cliente.
\anota{izq $= \rule{2cm}{0.4pt}$ \quad der $= \rule{2cm}{0.4pt}$ \quad $b = \rule{2cm}{0.4pt}$ mm}
\end{medicion}

\begin{medicion}{M5}{Desplazamiento longitudinal del marcador y altura}
\campo{Qué hacer} Del centro del marcador a la cara frontal de la placa delantera; y del
centro del marcador al extremo de las paletas. Aparte: \textbf{altura del marcador sobre
el piso}.

\campo{Qué predice el modelo} La configuración oficial declara altura
$= 90\ \mathrm{mm}$, sin confirmar. La distancia longitudinal no está en el modelo y hay
que medirla.

\campo{Qué significa} La altura entra en la corrección de paralaje
$\kappa = H/(H-h)$: si el marcador no está a $90\ \mathrm{mm}$, la posición reportada
tiene un error radial que crece hacia los bordes de la cancha. El desplazamiento
longitudinal define dónde está realmente el punto de contacto respecto de lo que
reporta la visión.
\anota{al frente $= \rule{2cm}{0.4pt}$ \quad a las paletas $= \rule{2cm}{0.4pt}$ \quad altura $= \rule{2cm}{0.4pt}$ mm}
\end{medicion}

\section*{Bloque C --- Lo que el DXF no podía dar}
\addcontentsline{toc}{section}{Bloque C}

\begin{medicion}{M6}{Alcance de las paletas $\lambda$}
\campo{Qué hacer} De la cara frontal de la placa delantera al extremo libre de las
paletas.

\campo{Qué predice el modelo} Nada: es el único parámetro que quedó sin valor. El modelo
lo esquiva con una cota, $\varrho_{\mathrm{cont}} \le 89{,}43\ \mathrm{mm}$.

\campo{Qué significa} Fija la distancia de contacto real y con ella el punto de puesta
$\varrho_{\mathrm{puesta}}$, que hoy vale $129{,}43\ \mathrm{mm}$ por cota. Un $\lambda$
grande acerca el punto de puesta y abarata cada maniobra.
\anota{$\lambda = \rule{4cm}{0.4pt}$ mm}
\end{medicion}

\begin{medicion}{M7}{Ancho \emph{real} máximo, incluyendo ruedas y motores}
\campo{Qué hacer} Lo más ancho del rover, punta a punta, sea lo que sea que sobresalga:
ruedas, ejes, tornillos, cables.

\campo{Qué predice el modelo} $W_{\mathrm{ext}} = 99{,}50\ \mathrm{mm}$, \textbf{bajo el
supuesto de que el chasis es la parte más ancha}.

\campo{Qué significa si es mayor} El radio envolvente $R_{\mathrm{rov}} = 68{,}44$ está
subestimado, y con él el inflado de obstáculos y la zona de exclusión de borde. Es un
error del lado peligroso: el planificador creería que pasa por huecos por los que no
pasa. \textbf{Esta es la medición más probable de refutar al modelo}, porque el archivo
de fabricación solo tiene las piezas de acrílico, no las ruedas.

\anota{$\rule{4cm}{0.4pt}$ mm \quad ¿qué es lo más ancho? $\rule{3cm}{0.4pt}$}
\end{medicion}

\begin{medicion}{M8}{Largo real máximo}
\campo{Qué hacer} Del punto más atrasado al más adelantado, incluyendo el ultrasónico,
las paletas y lo que sobresalga atrás.

\campo{Qué predice el modelo} El chasis mide $\Lambda = 94{,}00$, pero el largo total es
mayor: $94{,}00 + \lambda + $ lo que sobresalga atrás.

\campo{Qué significa} Igual que M7: entra en $R_{\mathrm{rov}}$ y en la condición de
voladizo sobre el precipicio.
\anota{$\rule{4cm}{0.4pt}$ mm}
\end{medicion}

\begin{medicion}{M9}{Paralelismo real de las paletas}
\campo{Qué hacer} Medir la separación entre caras interiores en \textbf{tres puntos}: en
la boca (extremo libre), a media altura, y en el fondo (contra la placa frontal).

\campo{Qué predice el modelo} Los tres valores iguales, $93{,}50\ \mathrm{mm}$.

\campo{Qué significa}
\begin{itemize}[leftmargin=1.6em,itemsep=2pt,topsep=3pt]
\item \textbf{Boca más ancha que el fondo} $\to$ hay algo de embudo. Buena noticia:
      mejora la captura, y el modelo pasaría a ser conservador en vez de exacto.
\item \textbf{Boca más angosta} $\to$ las paletas convergen. Mala noticia: el cubo puede
      entrar y trabarse.
\item \textbf{Iguales} $\to$ canal paralelo confirmado, el modelo es exacto.
\end{itemize}
\anota{boca $= \rule{2cm}{0.4pt}$ \quad medio $= \rule{2cm}{0.4pt}$ \quad fondo $= \rule{2cm}{0.4pt}$ mm}
\end{medicion}

\begin{medicion}{M10}{Altura de las paletas sobre el piso}
\campo{Qué hacer} Del piso al borde inferior de la paleta, y del piso al borde superior.

\campo{Qué predice el modelo} Nada: el modelo es plano y no mira alturas.

\campo{Qué significa} Si la paleta empuja por encima de la mitad del cubo
($30\ \mathrm{mm}$), el empuje genera un momento que tiende a \textbf{volcar} el cubo, y
la hipótesis H6 del modelo ---contacto plano, el cubo no vuelca--- deja de estar
garantizada. Si el borde inferior está muy alto, el cubo puede meterse por debajo.
\anota{inferior $= \rule{2.5cm}{0.4pt}$ \quad superior $= \rule{2.5cm}{0.4pt}$ mm}
\end{medicion}

\section*{Bloque D --- Los cubos}
\addcontentsline{toc}{section}{Bloque D}

\begin{medicion}{M11}{Arista real del cubo, medida por amplificación}
\campo{Qué hacer} Poné los \textbf{tres cubos en fila}, bien pegados, y medí el largo
total. Dividí por tres.

\campo{Qué predice el modelo} $3 \times 60 = 180{,}00\ \mathrm{mm}$.

\campo{Por qué así} Medir un cubo con regla da $\pm 0{,}5\ \mathrm{mm}$ de error; medir
tres y dividir da $\pm 0{,}17$. Es el mismo truco que se usó para verificar la escala
del archivo CAD.

\campo{Qué significa} $L$ entra en todo: en $w_S$, en $e_{\max}$, en el radio del
corredor. Un cubo de $62$ en vez de $60$ se lleva $1{,}4\ \mathrm{mm}$ más de la
tolerancia lateral en el peor caso.
\anota{$3L = \rule{3cm}{0.4pt}$ mm \quad$\Rightarrow$\quad $L = \rule{2cm}{0.4pt}$ mm}
\end{medicion}

\begin{medicion}{M12}{¿Las aristas del cubo están vivas o redondeadas?}
\campo{Qué hacer} Mirar y palpar una arista vertical del cubo. Si está redondeada,
estimar el radio comparándolo con algo conocido.

\campo{Qué predice el modelo} Aristas vivas: el modelo usa un cuadrado perfecto.

\campo{Qué significa} Un redondeo de radio $r$ reduce la anchura proyectada máxima de
$L\sqrt2$ a aproximadamente $L\sqrt2 - 2r(\sqrt2-1)$, lo que \textbf{devuelve} algo de
tolerancia: con $r = 2\ \mathrm{mm}$, $e_{\max}$ en el peor caso sube de $4{,}32$ a
$5{,}15\ \mathrm{mm}$ y el presupuesto deja de ser negativo por poco. También afecta al
detector de visión, que ajusta una silueta cuadrada.
\anota{vivas $\square$ \qquad redondeadas, $r \approx \rule{2cm}{0.4pt}$ mm}
\end{medicion}

\section*{Bloque E --- Para la capa de control (opcional)}
\addcontentsline{toc}{section}{Bloque E}

\begin{medicion}{M13}{Diámetro de rueda y trocha}
\campo{Qué hacer} Diámetro exterior de una rueda, y distancia entre los planos medios de
las dos ruedas (trocha).

\campo{Qué significa} Fuera del alcance del modelo geométrico, pero son los dos números
que convierten pulsos de encoder en milímetros y grados. Sin ellos, la odometría que
tiene que cubrir los huecos de oclusión no se puede calibrar. La trocha además fija el
radio mínimo de giro y con él la velocidad angular real $\omega$, que hoy se supone
$90\grados/\mathrm{s}$ en el costo de un quiebre.
\anota{diámetro $= \rule{2.5cm}{0.4pt}$ \quad trocha $= \rule{2.5cm}{0.4pt}$ mm}
\end{medicion}

\newpage
\section*{Resumen para llenar de un vistazo}
\addcontentsline{toc}{section}{Resumen}

\begin{center}
\small
\begin{tabular}{clrp{5.2cm}}
\toprule
\# & Qué & Predicho & Medido \\
\midrule
M1  & Hueco con cubo a $45\grados$ & $8{,}65$ & \\
M2  & Hueco con cubo alineado & $33{,}50$ & \\
--- & \emph{control:} M2 $-$ M1 & $24{,}85$ & \\
M3  & Lengüeta enrasada & sí & \\
\midrule
M4  & Marcador: izq / der & $49{,}75$ / $49{,}75$ & \\
M5  & Marcador: altura & $90$ & \\
M5  & Marcador: al frente / a paletas & --- & \\
\midrule
M6  & Alcance de paletas $\lambda$ & --- & \\
M7  & Ancho máximo real & $99{,}50$ & \\
M8  & Largo máximo real & $>94{,}00$ & \\
M9  & Canal: boca / medio / fondo & $93{,}50$ (los tres) & \\
M10 & Paletas: alto inferior / superior & --- & \\
\midrule
M11 & Tres cubos en fila & $180{,}00$ & \\
M12 & Aristas vivas & sí & \\
\midrule
M13 & Rueda: diámetro / trocha & --- & \\
\bottomrule
\end{tabular}
\end{center}

\vspace{0.5cm}

\begin{tcolorbox}[colback=black!5,colframe=black!65,boxrule=0.7pt,arc=1pt]
\textbf{Si solo podés hacer una medición, hacé M1.} Todo el resto del documento existe
para refinar; M1 decide si la conclusión central del trabajo se sostiene.

\smallskip
\textbf{Si podés hacer tres, hacé M1, M2 y M7.} M2 valida a M1, y M7 es la que tiene más
chance de encontrar un error del lado peligroso.
\end{tcolorbox}

\vspace{0.4cm}

\noindent
\footnotesize
Con los números anotados, el script \texttt{05\_Mediciones\_CAD/recalcular.py} rehace
todas las magnitudes derivadas y dice si algún resultado del informe cambia:

\begin{center}
\texttt{python3 recalcular.py --hueco45 8.7 --ancho-max 99.5 --lambda 40}
\end{center}

\end{document}
